% D.90. Root, FGH and PBC full proof/source/whole-certificate audits PASS.
\subsection{A nonunit five-valuation dominates the away lattice}

\begin{statement}{D.90}{Proved}{All compact elementary divisors when the five-valuation dominates}
Let
\[
 a=A\kappa(5)+\sum_{q\ne5}a_q\kappa(q),\qquad
 0\ne A\in\mathbb Z_5,\qquad a_q\longrightarrow0,
\]
and suppose the away-five class is nonzero. Put
\[
 \begin{aligned}
 m&=v_5(A),&r&=\min_{q\ne5}v_5(a_q),&0&\le m\le r,\\
 \lambda&=\sum_{q\ne5}a_q\frac{\log(q^4)}{4\log6},&
 N&=v_5(\lambda),&e&=N-r\in\mathbb Z_{\ge0}\cup\{\infty\}.
 \end{aligned}                                               \tag{D.90.1}
\]
For the actual cusp pair \(s=(a,0)\), the entire compact
kernel in D.63 is \(\mathcal K_s=\mathcal S_s\). Its compact
topological isomorphism type is
\[
 \mathcal S_s\simeq
 \begin{cases}
 0,&e=0,\\
 \displaystyle
 \left(\prod_{n\ge0}\mathbb Z/5^e\right)
 \times\prod_{j=1}^{\min(m,e-1)}(\mathbb Z/5^j)^2,
       &0<e<\infty,\\
 \displaystyle
 \left(\prod_{n\ge0}\mathbb Z_5\right)
 \times\prod_{j=1}^{m}(\mathbb Z/5^j)^2,
       &e=\infty.
 \end{cases}                                               \tag{D.90.2}
\]
Empty finite products are zero. The finite smaller-exponent
summands are retained. At \(m=0\) this recovers D.89.
The actual rational tangent \(630^5\) has
\((m,r,N,e)=(1,1,3,2)\), hence compact kernel
\[
             \left(\prod_{n\ge0}\mathbb Z/25\right)
                         \times(\mathbb Z/5)^2.        \tag{D.90.3}
\]
Every integer \(m\le r\) and positive finite integer \(e\)
is realized by the actual rational tangent
\(5^{5^m}(1+5^{e+1})^{5^r}\). The infinite endpoint is
realized by a finite-support completed cusp parameter.
Their actual metabelian fibres are nonempty and have
translation group \(\prod_{\mathbb N}\mathbb Z_5\), by D.74.

This is restricted to \(m\le r\). The intervening range
\(r<m<N\), arbitrary other dependent pairs, and the first
higher-derived obstruction are not settled here.
\end{statement}

\begin{proof}
Keep the actual notation \(\Lambda,u,v,t_1,t_2,w,P,I_0,N_0,N_1\)
of D.89. Set \(q_j=t_1^{5^j}-1\), including \(q_0=u\).
The actual common complement still acts as
\(\rho_A(g)=t_1^{A\kappa_5(g)}\rho_0(g)\).
Its two actual cyclotomic-kernel operators, one with
\(\kappa_5=0\) and primitive cubic coefficient and one
with \(\kappa_5=1\), give exactly
\[
       (\rho_A(g)-1:g\in N_0)=(P,t_1^A-1)=(P,q_m).
                                                               \tag{D.90.4}
\]
The proof is D.89's integral operator-ideal calculation before
its final use of \(A\) being a unit: if \(A=5^m\eta\),
then \((t_1^A-1)/q_m\) has unit constant coefficient
\(\eta\). Both generating elements are actual image elements.

\paragraph{The exact scalar module and its natural map.}
The actual global and local images are
\(5^r\mathbb Z_5\rtimes I_0\) and
\(5^N\mathbb Z_5\rtimes I_0\), with that same valuation
graph and \(5^\infty\mathbb Z_5=0\). If \(e=0\), they
agree and D.87 gives \(\mathcal S_s=0\).
Suppose \(e>0\). Put
\[
 R=\begin{cases}\Lambda/(q_N),&N<\infty,\\
                 \Lambda,&N=\infty,\end{cases}
                  \qquad J=q_rR.
\]
Every invariant functional on \(R\) annihilates \(q_mR\)
by (D.90.4). Since \(m\le r\), the actual inclusion
\(J\hookrightarrow R\) has image inside \(q_mR\).
Thus the natural invariant map is zero, exactly as a map,
and D.87 identifies
\[
 \mathcal S_s\simeq
              \bigl(J/(PJ+q_mJ)\bigr)(-1)_\Gamma,
                  \qquad \Gamma=\mathbb Z_5^\times.    \tag{D.90.5}
\]
In particular the whole finite-localization kernel is used;
no unrestricted everywhere-local kernel is substituted for it.

For finite \(N\), scalar division by the retained actual
generator \(q_r\) gives the relations
\(q_N/q_r,q_m,P\). Modulo \(q_m\), one has
\(t_1^{5^r}=1\), and hence the exact identity
\[
               q_N/q_r=\sum_{i=0}^{5^e-1}t_1^{i5^r}
                              \equiv5^e\pmod{q_m}.     \tag{D.90.6}
\]
At \(N=\infty\) there is no relative norm relation.
Consequently define
\[
 C=\begin{cases}\mathbb Z/5^e,&e<\infty,\\
                \mathbb Z_5,&e=\infty,\end{cases}
 \quad H_m=C[\mathbb Z/5^m],\quad
 B_C=H_m[[v]],\quad T_C=B_C/(uvw).                      \tag{D.90.7}
\]
Here \(H_0=C\). The compact module in (D.90.5) is
exactly \(T_C\) with the following actual action.
Substitution on \(q_r\) contributes
\(S_\chi(t_1^{5^r})=\chi\), which cancels the Cartier
untwist. The valuation graph multiplier is one modulo
\(q_m\). Therefore the action on the scalar module is
precisely the original \(\rho_0\), namely
\[
     S_\chi(t_1)S_\chi(t_2)\beta_g(t_1,t_2)
                            f(t_1^\chi,t_2^\chi).      \tag{D.90.8}
\]
It factors through \(\Gamma\) modulo \(P\).
This computation retains \(q_r\); treating it as a
weight-zero basis element would change the answer.

\paragraph{The exact torsion layer and zero snake map.}
Both \(v\) and \(w\) are nonzero divisors in \(B_C\).
For \(w\), the continuous change of variables
\(v\mapsto u+v+uv\), fixing \(u\), is an automorphism
of the completed ring and preserves the relation \(q_m\).
Thus it carries the injective multiplication by \(v\)
to multiplication by \(w\). In particular \(vw\)
is a nonzero divisor. We have an exact sequence
\[
 0\longrightarrow E_C\longrightarrow T_C\longrightarrow Q_C
                 \longrightarrow0,
 \quad E_C=vwB_C/(uvwB_C)\simeq C[[v]],
 \quad Q_C=B_C/(vw).                                    \tag{D.90.9}
\]
The action on \(E_C\), using its actual scalar generator
\(vw\), is \(\chi S_\chi(1+v)^3\) followed by
substitution. D.88's exact weighted-shell argument gives
\[
                         (E_C)_\Gamma\simeq\prod_{n\ge0}C.
                                                               \tag{D.90.10}
\]
For \(C=\mathbb Z_5\) this uses continuous functions valued
in \(\mathbb Q_5/\mathbb Z_5\); for finite \(C\) it uses
the finite coefficient dual. Both have the stated compact topology.

The quotient \(Q_C/uQ_C=C[[v]]/(v^2)\) has successive
tame characters \(\chi^2,\chi^3\). Thus every
\(\Gamma\)-invariant element of \(Q_C\) lies in \(uQ_C\).
The natural map \(uT_C\to uQ_C\) is an equivariant
isomorphism: modulo \(u\), the element \(vw\) is
\(v^2\), a nonzero divisor in \(C[[v]]\), so
\((u)\cap(vw)=(uvw)\). Every invariant lifts invariantly.
The snake lemma for an actual lift of the topological
generator \(2\in\Gamma\) therefore gives
\[
 0\longrightarrow(E_C)_\Gamma\longrightarrow(T_C)_\Gamma
                  \longrightarrow(Q_C)_\Gamma\longrightarrow0.
                                                               \tag{D.90.11}
\]
This exact injection remains valid over every \(C\) in (D.90.7).

\paragraph{The additional axis-intersection kernel is retained.}
The two-axis map for \(Q_C\) is not injective before tame
averaging, because \(u\) is a zero divisor in \(H_m\).
Write
\[
                    \mathcal N_m=\sum_{i=0}^{5^m-1}t_1^i.
\]
The literal cyclic coefficient equations give
\(\operatorname{Ann}_{H_m}(u)=C\mathcal N_m\).
Moreover \(w\mathcal N_m=v\mathcal N_m\).
If \(vf=wg\), then modulo \(v\),
\(u g(u,0)=0\), so \(g=c\mathcal N_m+vh\).
It follows that the kernel of the two-axis map is
generated by \(v\mathcal N_m\). This is an actual copy
of \(C\): if \(c v\mathcal N_m\) were in \((vw)\),
cancel \(v\) and evaluate on \(w=0\) to get
\(c\mathcal N_m=0\), hence \(c=0\).
There is consequently an exact sequence
\[
 0\longrightarrow C\,v\mathcal N_m\longrightarrow Q_C
 \longrightarrow H_m\oplus H_m\longrightarrow C\longrightarrow0.
                                                               \tag{D.90.12}
\]
The last map is the difference of augmentations.

The first term has exact tame character \(\chi^3\).
Indeed \(\sigma_\chi\mathcal N_m=\mathcal N_m\),
\(u\mathcal N_m=0\), and
\(v^2\mathcal N_m=vw\mathcal N_m=0\) in \(Q_C\).
Using (D.90.8) therefore gives
\(v\mathcal N_m\mapsto\chi^3 v\mathcal N_m\).
The final term has character \(\chi^2\). Both their
invariant and coinvariant parts under the order-four tame
subgroup vanish. Exact averaging removes these two terms
before taking the remaining procyclic coinvariants.

On the first boundary the action is
\(\chi S_\chi(t_1)\sigma_\chi\); on the diagonal it is
\(\chi S_\chi(t_1)t_1^{(1-\chi)/2}\sigma_\chi\).
The integral gauge \(t_1^{-1/2}\) identifies them, by D.39.
Multiplication by \(u\) transports the first action to
\(\chi\sigma_\chi\) on \(uH_m\).
Its kernel \(C\mathcal N_m\) has character \(\chi^2\);
the quotient \(H_m/uH_m\) for the target weighted action
has character \(\chi\). Exact tame averaging therefore
identifies each boundary's coinvariants with those of the
literal weighted permutation module
\[
                H_m,\qquad t^x\longmapsto\chi t^{\chi x}.
                                                               \tag{D.90.13}
\]

The zero orbit contributes nothing, since \(2-1\) is a unit.
The orbit of valuation \(i<m\) has length
\(4\cdot5^{m-i-1}\). Eliminating the basis vectors along
its actual weighted permutation cycle leaves the relation
\(2^{4\cdot5^{m-i-1}}-1\), whose five-valuation is
\(m-i\). Thus, naturally with respect to coefficient
reduction, (D.90.12)--(D.90.13) give
\[
                  (Q_C)_\Gamma\simeq
                            \bigoplus_{j=1}^{m}(C/5^jC)^2.
                                                               \tag{D.90.14}
\]
This integral calculation explains the smaller finite cyclic
factors; a rational rank computation would lose all of them.

\paragraph{Finite coefficients: split every finite-order lift.}
Let \(0<e<\infty\), so \(C=\mathbb Z/5^e\).
The entire module \((T_C)_\Gamma\) is killed by \(5^e\).
Choose a cyclic generator of order \(5^d\) in (D.90.14),
where \(d=\min(e,j)\), and any lift \(x\) to
\((T_C)_\Gamma\). Then \(5^d x\in(E_C)_\Gamma\)
and is killed by \(5^{e-d}\). In the actual product
\((E_C)_\Gamma\simeq\prod C\),
\[
      \operatorname{Ann}(5^{e-d})=5^d\prod C.
\]
Choose \(y\) in that product with \(5^dy=5^dx\).
Now \(x-y\) is a lift of the original cyclic generator
with order dividing \(5^d\). Doing this for the finitely
many independent cyclic generators produces a section of
(D.90.11). Its finite domain makes the section continuous;
the resulting compact-to-Hausdorff bijection is a topological
isomorphism. Hence
\[
 (T_C)_\Gamma\simeq
       \left(\prod_{n\ge0}\mathbb Z/5^e\right)
             \times\prod_{j=1}^{m}(\mathbb Z/5^{\min(e,j)})^2.
                                                               \tag{D.90.15}
\]
Only the factors with exponent exactly \(5^e\) can be
absorbed into the countable product. Retaining the others
gives the finite case of (D.90.2).

\paragraph{Infinite coefficients: the exact reduction maps split the extension.}
Now \(C=\mathbb Z_5\). Put
\(M=(T_{\mathbb Z_5})_\Gamma\),
\(E=(E_{\mathbb Z_5})_\Gamma\), and
\(Q=(Q_{\mathbb Z_5})_\Gamma\). Equation (D.90.11) is
\[
 0\longrightarrow E\simeq\prod\mathbb Z_5
       \longrightarrow M\longrightarrow
       Q\simeq\prod_{j=1}^{m}(\mathbb Z/5^j)^2
                                  \longrightarrow0.    \tag{D.90.16}
\]
It is not permissible to infer its splitting just from these
abstract end groups. We use the actual coefficient reductions.
For \(h\ge m\) (take \(h\ge1\) if \(m=0\)), reduction
of the displayed rings and operators gives
\[
 M/5^hM=(T_{\mathbb Z/5^h})_\Gamma,
 \qquad E/5^hE=(E_{\mathbb Z/5^h})_\Gamma.              \tag{D.90.17}
\]
These are the natural quotient maps: taking coinvariants is
the quotient by the image of the same actual operator minus
one and therefore commutes with coefficient reduction.
That image is compact, hence closed. The explicit cycle
calculation (D.90.14) identifies the corresponding map on
\(Q\) with \(Q\to Q/5^hQ=Q\).

The zero-snake injection at the finite coefficient level in
(D.90.11) implies
\[
                         E\cap5^hM=5^hE.
\]
Since \(5^hQ=0\), one also has \(5^hM\subset E\).
Consequently
\[
                              5^hM=5^hE.               \tag{D.90.18}
\]
The finite-coefficient splitting already proved gives an
order-\(5^j\) lift of each cyclic generator of \(Q\)
in \(M/5^hM\). Lift it to \(x\in M\). Then
\(5^j x\in5^hM=5^hE\subset5^jE\), because \(j\le m\le h\).
An adjustment by an element of \(E\) makes \(5^j x=0\)
without changing its image in \(Q\). Choosing the finitely
many such lifts gives an actual continuous section of
(D.90.16). This proves its splitting and the infinite case
of (D.90.2). All maps here are the actual coefficient
reductions; neither matching cardinalities nor an unproved
interchange of limits is used.

\paragraph{The complete compact kernel and actual examples.}
Both the global depth and the full local depth of \(a\)
are \(m\): its away depth is \(r\ge m\), its unit local
depth is \(N\ge r\), and its five-valuation depth is
\(m\). The same holds for \(a-b=a\). The zero boundary
\(b\) contributes an identical countable product on both
sides. D.77's natural boundary map is therefore an
isomorphism, and D.63 gives \(\mathcal K_s=\mathcal S_s\).
This also completes the zero-index case.

For the stated finite examples take \(A=5^m\) and the
away class \(5^r\kappa(1+5^{e+1})\). Its valuation at
two is exactly \(5^r\), and its local unit depth is
\(r+e\), by the proved D.88 logarithm calculation.
For \(m=r=1,e=2\), the tangent is
\(5^5\cdot126^5=630^5\), giving (D.90.3).
For infinite index take
\[
 a=5^m\kappa(5)+
    5^r\bigl(\lambda_3\kappa(2)-\lambda_2\kappa(3)\bigr),
 \qquad \lambda_q=\log(q^4)/(4\log6).
\]
The away class is nonzero of depth \(r\) and has zero
local coefficient. These are actual full cusp parameters;
their finite valuation support permits the D.74 fibre
calculation. No assertion about all other abelian fibres follows.

The literal integral boundary matrices, exact norm-line and
half-gauge identities, and Smith forms through cyclic depth
three are retained in \path{code/de_sha_dominant_valuation_cusp.py}
and its executed outputs. These finite checks support the
actual formulas; the extension and infinite-endpoint proofs
above supply the all-depth conclusion.
\end{proof}
